\section{Metodología de búsqueda}
\label{sec:metodologia}

La búsqueda se realizó en septiembre de 2026 mediante Consensus para localizar literatura académica y búsquedas web para identificar documentos técnicos y casos recientes de adopción. En Consensus se utilizaron las cadenas:

\begin{itemize}
    \item ``post-quantum cryptography'' + review, standardization, migration.
    \item ``post-quantum cryptography'' + implementation, performance.
    \item ``post-quantum cryptography'' + NIST, migration, security.
\end{itemize}

También se utilizó Gemini como apoyo para identificar posibles fuentes y casos de estudio, pero las recomendaciones fueron verificadas posteriormente en las fuentes originales.

\subsection{Preguntas guía}

La revisión se orientó a responder: ¿qué riesgos introduce la computación cuántica?, ¿qué soluciones y estándares existen?, ¿qué dificultades técnicas presenta la migración? y ¿qué evidencia existe de adopción actual?

\subsection{Criterios de selección}

Se seleccionaron cinco fuentes combinando las tres categorías permitidas: académicas, técnicas y de actualidad, con un mínimo de dos académicas. Se priorizaron fuentes con autoría identificable, relación directa con criptografía post-cuántica, evidencia verificable y aporte complementario al resto de trabajos. Se excluyeron fuentes sin autoría clara, contenido secundario cuando existía una fuente primaria y trabajos poco relacionados con las preguntas guía.

\subsection{Alcance y limitaciones}

La búsqueda se concentró principalmente en fuentes recientes en inglés y español. Debido al volumen de literatura, se realizó un filtrado inicial mediante título, resumen, introducción y conclusiones, seguido de una revisión más detallada de las fuentes seleccionadas. Consensus no necesariamente ordena los trabajos según su utilidad para este caso particular, y pueden existir trabajos relevantes en otros idiomas o con menor visibilidad que no hayan sido identificados.